\chapter*{About These Notes}
\addcontentsline{toc}{chapter}{About These Notes}
These notes combine the queueing primer and the disaggregated-serving course.
They assume real analysis and mathematical statistics, but no prior queueing theory.
Definitions, proofs and exercises for the common tools appear once. The final two
chapters apply them to a colocated replica and a prefill--decode deployment.

\paragraph{The system.} Prefill reads a prompt and produces the first token.
Decode produces the remaining tokens. Both need KV memory. A session returns
with another turn after a tool call, so session arrivals and turn arrivals are
different processes. A cache policy changes how much prefill work a turn needs.
A colocated engine shares iterations between prefill and decode; a disaggregated
engine transfers KV between their separate memory pools.

\paragraph{Reading plan.} \Cref{sec:L1} introduces both architectures, serQ v0.1.0,
Poisson arrivals, PASTA, renewal rewards and Little's law. \Cref{sec:L2} develops
Markov chains, M/M/1, memory slots, processor sharing and insensitivity.
\Cref{sec:L3} proves Pollaczek--Khinchine and the price of a miss.
\Cref{sec:L4} covers finite populations, resident memory and eviction.
\Cref{sec:L5} builds the colocated scheduler and its holding-time rule.
\Cref{sec:L6} adds transfer prices, two memory pools and cache feedback.
The final appendix records what a release comparison can and cannot establish.

\paragraph{Language and models.} Executable examples use serQ v0.1.0 (IR v9).
Mathematical queues remain idealisations with explicit assumptions; a step engine
with block admission and preemption does not inherit their formulas automatically.
The original paper concerns colocated serving. Disaggregation is a course extension.
The exercises reinforce the proofs; \cite{kleinrock1975,harchol2013} are introductory
references and \cite{asmussen2003} gives a rigorous treatment.
